\section{DADOS E MÉTODO}

Este capítulo apresenta a base, delimita o que ela permite medir e descreve como o texto das avaliações é convertido em medidas por aspecto. A validação dessas medidas é tratada na Seção 4, e as análises que as utilizam, nas Seções 5 a 7.

\subsection{Base e recorte}

A base provém do projeto Inside Airbnb \citep{insideairbnb} e reúne 42.013 anúncios do Rio de Janeiro e 995.072 avaliações textuais. Entre os anúncios, 33.140 possuem nota global e as seis subnotas oficiais (limpeza, localização, check-in, comunicação, custo-benefício e precisão do anúncio). Foram considerados elegíveis os que têm nota maior que zero e ao menos duas avaliações, o que resulta em 28.890 anúncios.

O custo computacional da inferência levou a um recorte: o sentimento foi calculado sobre as 30 avaliações mais recentes de cada anúncio. O recorte contém 475.858 avaliações, depois da remoção de uma avaliação automática e de 311 duplicatas, e o limite atinge 32\% dos anúncios elegíveis (9.266 anúncios têm mais avaliações do que o teto). As avaliações foram divididas em sentenças, o que produziu 1.512.063 unidades, das quais 1.366.091 são textos únicos, com média de 3,18 sentenças por avaliação e 12,4 \textit{tokens} por sentença. O idioma foi detectado por avaliação: português (66,2\%), inglês (15,5\%), espanhol (14,2\%), francês (2,4\%), alemão (1,1\%) e italiano (0,3\%). O corpus não foi separado por idioma, e os dois modelos de sentimento adotados são multilíngues.

A nota oficial é fortemente comprimida no topo da escala. Entre os 28.890 anúncios elegíveis, 9.373 (32,4\%) têm nota exatamente igual a 5,0.

\subsection{O que a base permite medir}

Os dados não trazem ocupação nem receita observadas. Os campos \texttt{estimated\_occupancy\_l365d} e \texttt{estimated\_revenue\_l365d} são estimativas calculadas pelo Inside Airbnb e, como a verificação deste estudo mostrou, correspondem a funções exatas de variáveis do próprio anúncio. A ocupação estimada é dada por $\min(255,\; R_{12} \times 2 \times \max(3, N_{\min}))$, em que $R_{12}$ é o número de avaliações nos últimos 12 meses e $N_{\min}$ é a estadia mínima. A fórmula reproduz o campo em 100,0\% dos 42.013 anúncios. A receita estimada é igual ao preço multiplicado pela ocupação estimada em 100,0\% dos 38.356 anúncios com preço informado, e a ocupação é zero nos 12.979 anúncios sem avaliação nos últimos 12 meses (100,0\% dos casos). A construção segue o modelo descrito pelo Inside Airbnb, que converte avaliações em reservas com taxa de avaliação de 50\% e estadia média de três noites \citep{insideairbnb}. Na faixa de nota 5,0, a correlação de postos entre a ocupação estimada e o número de avaliações dos últimos 12 meses é 0,97.

Por esse motivo, os dois campos são tratados como índices de volume recente de avaliações, e os termos ocupação e receita só aparecem acompanhados dessa qualificação. Há uma consequência de desenho. Controlar uma regressão pelo número de avaliações dos últimos 12 meses equivale a controlar o próprio desfecho, e por isso esse controle não é utilizado. O objeto central do estudo são a nota global, as subnotas e o texto das avaliações. O volume de avaliações entra como desfecho auxiliar em três formas: o índice de ocupação, o índice de receita (em logaritmo) e o número de avaliações dos últimos 12 meses (em logaritmo), este último sem construção algébrica.

\subsection{Medição do sentimento}

O sentimento de cada sentença é medido por dois modelos, porque a escolha de um único modelo poderia alterar as conclusões. O modelo A (\texttt{cardiffnlp/twitter-xlm-roberta-base-sentiment}) é um classificador multilíngue treinado em mensagens de uma rede social \citep{barbieri2022,hf_cardiffnlp}, com três classes: negativo, neutro e positivo. Sua saída foi convertida em estrelas esperadas pela expressão $1 \cdot P(\text{neg}) + 3 \cdot P(\text{neu}) + 5 \cdot P(\text{pos})$. O modelo B (\texttt{nlptown/bert-base-multilingual-uncased-sentiment}) é um modelo multilíngue da família BERT \citep{devlin2019} ajustado a avaliações de produtos \citep{hf_nlptown}, com cinco classes que correspondem a 1 a 5 estrelas. As estrelas esperadas são $\sum_k k \cdot P(k)$.

A polaridade discreta de A é a mais provável entre as três classes. Em B, a polaridade é negativa quando a classe mais provável tem 1 ou 2 estrelas, neutra quando tem 3 e positiva quando tem 4 ou 5. Os modelos foram aplicados às mesmas 1.512.063 sentenças (1.366.091 textos únicos), com a biblioteca Transformers \citep{wolf2020}. As estrelas médias por texto único são 4,19 em A e 4,06 em B, e a proporção de sentenças com quatro estrelas ou mais é 75,1\% e 70,0\%. As polaridades (negativa, neutra e positiva) distribuem-se em 6,8\%, 14,3\% e 78,9\% em A, e em 10,4\%, 8,1\% e 81,5\% em B. Nos dois casos a maior parte das sentenças é positiva, o que reproduz no texto o viés de positividade das avaliações online \citep{zervas2021,bridges2018}.

As análises seguintes usam a ordenação do sentimento, e não o seu nível absoluto. Importa, portanto, que as estrelas esperadas ordenem as sentenças como um leitor humano. A correlação de postos entre o rótulo manual ($-1$, 0, $+1$) e as estrelas esperadas é 0,64 em A e 0,63 em B.

Antes de executar o modelo B sobre o corpus, fixou-se a regra de escolha do modelo principal. O modelo B seria adotado como principal se o intervalo de 95\% da diferença B menos A no kappa ponderado linear populacional, estimado por \textit{bootstrap} estratificado \citep{efron1979}, excluísse zero a favor de B. Caso contrário, o principal seria A. O outro modelo é reportado em todos os resultados como análise de robustez. O kappa populacional é calculado sobre sentenças rotuladas manualmente e ponderadas pelo tamanho dos estratos de origem, como detalhado na Seção 4.

\subsection{Tópicos e aspectos}

As sentenças foram agrupadas pelo BERTopic \citep{grootendorst2022}, que combina \textit{embeddings} de sentenças, redução de dimensão e agrupamento por densidade. Foram usados os \textit{embeddings} multilíngues \texttt{paraphrase-multilingual-MiniLM-L12-v2} \citep{reimers2020}, a redução por UMAP \citep{mcinnes2018umap} e o agrupamento por HDBSCAN \citep{mcinnes2017}. O ajuste foi feito em uma amostra de 200.000 sentenças, e o modelo foi aplicado às demais. Resultaram 185 tópicos. Na amostra de ajuste, 29,8\% das sentenças foram classificadas como ruído; depois da reatribuição ao centroide mais próximo, o ruído caiu a 1,0\% dos textos. Distinguem-se, assim, as sentenças com tópico direto (atribuídas pelo próprio BERTopic) e as reatribuídas, cuja atribuição é menos precisa (Seção 4). A coerência $C_v$ \citep{roder2015} média dos tópicos é 0,388.

Como método de checagem, ajustou-se também um modelo LDA \citep{blei2003} com $K = 10$, valor que maximiza o $C_v$ em uma grade de 5 a 40 tópicos e produz $C_v$ médio de 0,625. Os valores de $C_v$ descrevem cada método e não os colocam em disputa. O LDA cumpre apenas o papel de verificar se os resultados por aspecto dependem do método de tópicos (Seção 4).

Cada tópico do BERTopic foi associado a um de dez aspectos: limpeza, localização, check-in, comunicação, custo-benefício, ruído e conforto, segurança, precisão do anúncio, comodidades e satisfação geral (este último reúne elogios e queixas genéricos). O código sugeriu o aspecto pela similaridade do tópico com termos-semente multilíngues, e a decisão final foi tomada por um pesquisador, tópico a tópico. Dos 185 tópicos, 105 mantiveram a sugestão automática e 80 foram alterados na revisão. A satisfação geral concentra 53 tópicos. Cada sentença herda o aspecto do tópico a que pertence.

\subsection{Perfil por anúncio}

Para cada anúncio e cada aspecto, o perfil é a média do sentimento das sentenças do aspecto, calculada quando há ao menos duas sentenças. Ao todo, 28.889 anúncios têm perfil. A Tabela~\ref{tab:aspectos} mostra a participação de cada aspecto no corpus e o número de anúncios para os quais o perfil está disponível.

\begin{table}[htb]
\centering
\caption{Participação dos aspectos nas sentenças e anúncios com perfil disponível.}
\small
\begin{tabular}{l|r|r}
\hline
Aspecto & Sentenças do corpus & Anúncios com o aspecto\\
\hline
Satisfação geral & 27,1\% & 26.378\\
Comunicação & 19,2\% & 25.418\\
Localização & 18,3\% & 24.101\\
Comodidades & 13,0\% & 22.272\\
Limpeza & 7,6\% & 19.095\\
Precisão do anúncio & 5,2\% & 16.638\\
Ruído e conforto & 2,5\% & 9.550\\
Check-in & 2,3\% & 8.919\\
Segurança & 1,3\% & 5.072\\
Custo-benefício & 1,0\% & 3.552\\
Fora do perfil (\texttt{geral}) & 2,5\% & --\\
\hline
\end{tabular}
\label{tab:aspectos}
\end{table}

\subsection{Desenho analítico}

Todas as análises estimam associações e não efeitos causais. O alvo principal é o percentil da nota global, escolhido porque a distribuição da nota se concentra no topo. A importância dos aspectos é decomposta pelo método LMG, equivalente à média de Shapley do $R^2$ sobre todas as ordens de entrada dos preditores \citep{lindeman1980,gromping2006}. A decomposição é complementada por um modelo de árvores com gradiente (GBM), que aceita valores ausentes, com importância por permutação fora da amostra, por uma regressão logística para a nota igual ou superior a 4,9 e por coeficientes padronizados com erro-padrão robusto HC3. Como verificação de sensibilidade, a decomposição é repetida com os tópicos diretos, com os anúncios de até 30 avaliações, com ao menos três aspectos observados e sem a satisfação geral entre os preditores, e por segmento (tipo de acomodação, idioma predominante e os maiores bairros).

A leitura por tópico compara o percentil da nota entre anúncios com e sem ao menos uma sentença negativa (ou positiva) no tópico, com controle do volume de avaliações e do número de sentenças. Os testes de tópicos usam a correção de Benjamini e Hochberg \citep{benjamini1995} para comparações múltiplas.

O índice sintético é a média simples do sentimento dos aspectos disponíveis do anúncio, em estrelas de 1 a 5, com exigência de ao menos três aspectos específicos com evidência; a satisfação geral entra na média, mas não conta para o mínimo. Seu poder explicativo sobre os índices de volume é avaliado por modelos aninhados, com controles de volume total de avaliações (em logaritmo), preço (em logaritmo), tipo de acomodação e bairro. A validação é feita por validação cruzada com cinco dobras repetidas cinco vezes, com a mesma semente e as mesmas dobras em todos os modelos \citep{pedregosa2011}.

Todos os números apresentados nos resultados foram extraídos de uma planilha única gerada pelo \textit{pipeline} de análise, com carimbo de data e hora, e conferidos um a um contra o texto.
